% TODO list
% Explain why verifier polytope cannnot be directly used to generate safe mppi data 
% implicit safe set necessary for the real-time, highly parallel sample rejection step of MPPI. In contrast, the verifier polytope requires solving a trajectory-dependent optimization problem, making it computationally prohibitive for real-time MPPI sampling, though highly effective for offline or lower-frequency evaluation of the generative model's outputs.
% No implicit safeset 
% Titles on the y axis in Fig. 1 
% Eq (14) fix
% Nominal polytope & Verifier polytope 의 motivation? Easy-to-simulate. 
% Gamma 의 능력을 좀 더 부각시킬 것 (Figure 4) 
%   SFM 을 그대로 할 지 모르겠음 Fix 된 pedestrian dataset 에 대해 random seed 로 coverage 측정
% Self generation 을 offline simulation 상에서 한다고 말해야 함
% 	okay
% Coverage 가 늘어가는걸 보고 싶다? 후속 연구에서. 
% 	okay 3d 에 대해 하고 있음
% Novelty 에 대한 정의가 잘 정의되어야 한다. 다양한 valid behavior 가 나오는지가 궁금하다 
% 	verifier uncertainty
% Validity -> 100% 가 되는 guarantee. 
% 	TBD
% Argmax term in algorithm 을 언급해보자
% 	다양한 option 이 있을 수 있고, 이를 Figure 3 에서 실험하면 좋을 듯
% Verifier 가 obstacle avoidance 아니고 humanoid 가 성공적으로 수행할 확률
% 	후속 연구로 좋음
% Adjacency graph 랑 combined 된 global knowledged generative model
%   후속 연구
% Define the OOD concretely (e.g., obstacles)
\note[Dohyun]{If we present the system as a reduced-order model, the real robot is still governed by a more complicated full-order system. So we need to be clear about the modeling gap. Later safety claims depend on how well the reduced-order trajectory can be tracked by the full-order system. We could add this uncertainty later, but it should be acknowledged.}

\note{Structure for each subsection: 1. Define data used for training, In-distribution and OOD, 2. Safe Flow Expansion, how we subsequently gather positive samples (execution strategy) }

\note{Proposing gamma scheduling now can handle the performance side.}

